\documentclass[11pt]{article}
\usepackage{notesstyle}

\begin{document}
\notesheader{3.5}{Hierarchical Indexing}{The \code{MultiIndex} --- packing extra dimensions of labels into a 1-D \code{Series} or 2-D \code{DataFrame}.}

\section{Why you would ever want a second level of index}

So far an index has been a flat list of labels: one label per row. That is enough for genuinely two-dimensional data (rows $\times$ columns). But an enormous amount of real data is naturally \emph{three} or more dimensions: sales by \emph{store} and \emph{month}, gene expression by \emph{patient} and \emph{tissue}, sensor readings by \emph{location} and \emph{timestamp}. You want to keep indexing data --- looking it up by meaningful labels --- but now along more than one axis at once.

\begin{background}
The cleanest tool for $N$-dimensional labeled data is a genuinely $N$-dimensional container (the \code{xarray} library exists for exactly this). But in day-to-day work we usually want to stay inside the \code{Series}/\code{DataFrame} world, because all the alignment, grouping, and plotting machinery lives there. \textbf{Hierarchical indexing} (a.k.a. multi-level indexing) is Pandas's answer: it encodes the extra axes as nested \emph{levels} of a single index, so a 1-D \code{Series} can stand in for 2-D data, and a 2-D \code{DataFrame} can stand in for 3-D or 4-D data.
\end{background}

\subsection{The naive approach: Python tuples as keys}

Suppose we have a population figure for a few US states, but at \emph{two different years}. The label for each value is really a pair \code{(state, year)}. A first instinct is to use tuples as the index keys:

\begin{lstlisting}
import numpy as np
import pandas as pd

index = [('California', 2010), ('California', 2020),
         ('New York',   2010), ('New York',   2020),
         ('Texas',      2010), ('Texas',      2020)]
populations = [37_253_956, 39_538_223,
               19_378_102, 20_201_249,
               25_145_561, 29_145_505]
pop = pd.Series(populations, index=index)
print(pop)
\end{lstlisting}

\begin{outblock}
(California, 2010)    37253956
(California, 2020)    39538223
(New York, 2010)      19378102
(New York, 2020)      20201249
(Texas, 2010)         25145561
(Texas, 2020)         29145505
dtype: int64
\end{outblock}

This \emph{works}: \code{pop[('California', 2020)]} returns the right number. But it is clumsy the moment you want a \emph{slice}. The index is, as far as Pandas is concerned, a flat list of opaque tuple objects. There is no notion that the first element is a ``state level'' and the second a ``year level.'' So a natural question like ``give me every state's 2020 figure'' has no clean expression:

\begin{lstlisting}
# Select all years 2020 -- ugly: we must loop the tuples by hand
pop[[i for i in pop.index if i[1] == 2020]]
\end{lstlisting}

\begin{outblock}
(California, 2020)    39538223
(New York, 2020)      20201249
(Texas, 2020)         29145505
dtype: int64
\end{outblock}

That list comprehension is exactly the kind of manual bookkeeping Pandas is supposed to eliminate. It is also slow (a Python loop over every label) and it does not generalize: slicing on the \emph{outer} level, partial selection, sorting by one level --- all become bespoke code.

\begin{keyidea}
A list of tuples \emph{looks} hierarchical but is not. Pandas sees a flat sequence of tuple objects with no level structure, so none of the level-aware operations (partial indexing, cross-section, group-by-level) are available. The fix is to promote those tuples into a real \code{MultiIndex}.
\end{keyidea}

\subsection{The clean solution: \code{pd.MultiIndex}}

A \code{MultiIndex} is an index object that explicitly stores \emph{multiple levels}, each its own array of labels, plus the integer codes that say which label of each level a given row uses. We can build one straight from our tuples:

\begin{lstlisting}
mi = pd.MultiIndex.from_tuples(index)
pop = pop.reindex(mi)
print(pop)
\end{lstlisting}

\begin{outblock}
California  2010    37253956
            2020    39538223
New York    2010    19378102
            2020    20201249
Texas       2010    25145561
            2020    29145505
dtype: int64
\end{outblock}

Two things changed. Visually, the repeated outer label (\code{California}, \code{New York}, \dots) is now collapsed --- the blank space under \code{California} means ``same as above,'' signalling the level structure. Functionally, we now get the slice we wanted with trivial syntax:

\begin{lstlisting}
pop[:, 2020]      # all rows whose SECOND level == 2020
\end{lstlisting}

\begin{outblock}
California    39538223
New York      20201249
Texas         29145505
dtype: int64
\end{outblock}

The result is a single-level \code{Series} indexed only by state --- Pandas dropped the level we fully specified. This is shorter, faster, and far more readable than the list comprehension.

\section{A MultiIndex is literally an extra dimension}

Here is the mental model worth internalizing. Our \code{pop} Series has six entries laid out along one axis, but its \emph{meaning} is a $3 \times 2$ grid (three states, two years). The \code{MultiIndex} is what lets a 1-D container carry 2-D meaning. The conversion between ``stacked into a column'' and ``spread into a grid'' is so common it has two named methods: \code{unstack} and \code{stack}.

\begin{center}
\begin{tikzpicture}[font=\sffamily\small,
   lvl/.style={font=\sffamily\footnotesize\color{accent}},
   cell/.style={draw=rulegray, minimum width=2.0cm, minimum height=0.62cm, anchor=center},
   outer/.style={cell, fill=bgblue, font=\sffamily\small\color{accent}},
   inner/.style={cell, fill=bggreen, font=\sffamily\small\color{accent2}},
   val/.style={cell, fill=white},
   col/.style={cell, fill=bggreen, font=\sffamily\small\color{accent2}, minimum width=1.6cm}]

  % ---- LEFT: multiply-indexed Series ----
  \node[lvl] at (1.0,4.55) {\textbf{Series} with 2-level index};
  % California
  \node[outer, minimum height=1.24cm] (ca) at (0,3.5) {California};
  \node[inner] (ca1) at (1.5,3.81) {2010};
  \node[inner] (ca2) at (1.5,3.19) {2020};
  \node[val]   at (3.4,3.81) {37253956};
  \node[val]   at (3.4,3.19) {39538223};
  % New York
  \node[outer, minimum height=1.24cm] (ny) at (0,2.26) {New York};
  \node[inner] at (1.5,2.57) {2010};
  \node[inner] at (1.5,1.95) {2020};
  \node[val]   at (3.4,2.57) {19378102};
  \node[val]   at (3.4,1.95) {20201249};
  % Texas
  \node[outer, minimum height=1.24cm] (tx) at (0,1.02) {Texas};
  \node[inner] at (1.5,1.33) {2010};
  \node[inner] at (1.5,0.71) {2020};
  \node[val]   at (3.4,1.33) {25145561};
  \node[val]   at (3.4,0.71) {29145505};

  \node[lvl] at (0,0.05)   {outer level};
  \node[lvl] at (1.5,0.05) {inner level};
  \node[lvl] at (3.4,0.05) {values};

  % ---- arrows ----
  \draw[-{Stealth[length=3mm]}, accent, line width=1pt]
      (4.55,2.7) -- node[above, font=\sffamily\footnotesize\color{accent}]{unstack()}
                    node[below, font=\sffamily\footnotesize\color{accent}]{(inner $\to$ cols)} (6.6,2.7);
  \draw[-{Stealth[length=3mm]}, accent2, line width=1pt]
      (6.6,2.15) -- node[above, font=\sffamily\footnotesize\color{accent2}]{stack()} (4.55,2.15);

  % ---- RIGHT: 2-D DataFrame ----
  \node[lvl] at (8.7,4.55) {\textbf{DataFrame} (2-D grid)};
  \node[col] at (8.1,3.95) {2010};
  \node[col] at (9.7,3.95) {2020};
  \node[outer, minimum width=1.8cm] at (6.55,3.33) {California};
  \node[outer, minimum width=1.8cm] at (6.55,2.71) {New York};
  \node[outer, minimum width=1.8cm] at (6.55,2.09) {Texas};
  \node[val, minimum width=1.6cm] at (8.1,3.33) {37253956};
  \node[val, minimum width=1.6cm] at (9.7,3.33) {39538223};
  \node[val, minimum width=1.6cm] at (8.1,2.71) {19378102};
  \node[val, minimum width=1.6cm] at (9.7,2.71) {20201249};
  \node[val, minimum width=1.6cm] at (8.1,2.09) {25145561};
  \node[val, minimum width=1.6cm] at (9.7,2.09) {29145505};
\end{tikzpicture}
\end{center}

\begin{lstlisting}
pop_df = pop.unstack()   # inner level (year) becomes the columns
pop_df
\end{lstlisting}

\begin{center}
\renewcommand{\arraystretch}{1.2}
\begin{tabular}{l r r}
\toprule
 & \textbf{2010} & \textbf{2020} \\
\midrule
California & 37253956 & 39538223 \\
New York   & 19378102 & 20201249 \\
Texas      & 25145561 & 29145505 \\
\bottomrule
\end{tabular}
\end{center}

\begin{lstlisting}
pop_df.stack()   # columns fold back into a second index level
\end{lstlisting}

And \code{stack} returns us exactly to the original multiply-indexed \code{Series}. The pair is a lossless round trip.

\begin{intuition}
\code{unstack} pivots one index level ``up'' to become columns (data gets \emph{wider}); \code{stack} pivots columns ``down'' to become an index level (data gets \emph{taller}). It is the same numbers, reshaped --- a \code{MultiIndex} on rows and a column axis are two views of the same higher-dimensional cube.
\end{intuition}

\subsection{Why bother --- just add a column?}

Why not keep a plain \code{DataFrame} with a \code{year} \emph{column} instead of a year index level? You can, and often should. But once a quantity is a label rather than a value, treating it as an index level unlocks alignment and group-by-level operations, lets you add \emph{more} columns of actual measurements without confusion, and gives you the \code{stack}/\code{unstack} reshaping for free. The \code{MultiIndex} earns its keep precisely when you have several measured columns indexed by the same set of label axes:

\begin{lstlisting}
pop_df = pd.DataFrame({'total': pop,
                       'under18': [9_284_094, 8_898_092,
                                   4_687_374, 4_318_033,
                                   6_879_014, 7_503_437]})
pop_df['frac_under18'] = pop_df['under18'] / pop_df['total']
pop_df['frac_under18'].unstack()
\end{lstlisting}

\begin{outblock}
                2010      2020
California   0.249265  0.225041
New York     0.241886  0.213750
Texas        0.273568  0.257619
dtype: float64
\end{outblock}

The arithmetic ran row-by-row honoring both index levels, then \code{unstack} reshaped the answer into a readable grid --- no manual joining of state and year keys anywhere.

\section{Constructing a MultiIndex}

You rarely call \code{from\_tuples} by hand. More often a \code{MultiIndex} appears \emph{implicitly}, or you build it from arrays/products.

\subsection{Implicit construction}

If you pass a \emph{list of two or more lists} as the \code{index} to a \code{Series} or \code{DataFrame}, Pandas builds the \code{MultiIndex} for you:

\begin{lstlisting}
df = pd.DataFrame(np.round(np.random.rand(4, 2), 2),
                  index=[['a', 'a', 'b', 'b'], [1, 2, 1, 2]],
                  columns=['data1', 'data2'])
\end{lstlisting}

Likewise, passing a dictionary whose \emph{keys are tuples} makes Pandas recognize the levels automatically:

\begin{lstlisting}
data = {('California', 2010): 37_253_956,
        ('California', 2020): 39_538_223,
        ('Texas',      2010): 25_145_561,
        ('Texas',      2020): 29_145_505}
pd.Series(data)   # -> a MultiIndexed Series, no extra work
\end{lstlisting}

\subsection{Explicit constructors}

When you want to be deliberate, the \code{pd.MultiIndex} class has three workhorse class methods:

\begin{lstlisting}
# From a list of (level0, level1) tuples
pd.MultiIndex.from_tuples([('a', 1), ('a', 2), ('b', 1), ('b', 2)])

# From a list of arrays, one array per level (parallel, row-aligned)
pd.MultiIndex.from_arrays([['a', 'a', 'b', 'b'], [1, 2, 1, 2]])

# From the CARTESIAN PRODUCT of single-level iterables
pd.MultiIndex.from_product([['a', 'b'], [1, 2]])
\end{lstlisting}

\begin{intuition}
Choose by how your levels relate. \code{from\_product} when every combination of level values exists (a full grid --- e.g. every state crossed with every year). \code{from\_arrays} when you already have the levels as parallel columns and not all combinations occur. \code{from\_tuples} when the pairs arrive naturally as tuples. They all produce the same kind of object.
\end{intuition}

\subsection{Naming the levels}

An unnamed level is hard to refer to later. Every \code{MultiIndex} has a \code{names} attribute --- a list with one name per level --- and you should set it:

\begin{lstlisting}
pop.index.names = ['state', 'year']
print(pop)
\end{lstlisting}

\begin{outblock}
state       year
California  2010    37253956
            2020    39538223
New York    2010    19378102
            2020    20201249
Texas       2010    25145561
            2020    29145505
dtype: int64
\end{outblock}

Now the level names print as a header, and you can use them in \code{groupby(level='year')}, \code{xs(level='state')}, and so on, instead of remembering that ``year is level 1.''

\section{Hierarchical indexing on columns too}

Rows are not special. A \code{DataFrame}'s \emph{columns} are themselves an \code{Index}, so they can be a \code{MultiIndex} as well. This gives you fully four-dimensional tables: two label axes down the rows, two across the columns. A classic use is repeated measurements --- several subjects, each measured for several quantities, across several visits.

\begin{lstlisting}
# hierarchical row index: (year, visit)
rows = pd.MultiIndex.from_product([[2024, 2025], [1, 2]],
                                  names=['year', 'visit'])
# hierarchical column index: (subject, measurement-type)
cols = pd.MultiIndex.from_product([['Ada', 'Bo'], ['hr', 'temp']],
                                  names=['subject', 'type'])

rng = np.random.default_rng(0)
data = np.round(37 + 3 * rng.standard_normal((4, 4)), 1)
data[:, ::2] = np.round(data[:, ::2] * 2)   # make hr look heart-rate-ish
health = pd.DataFrame(data, index=rows, columns=cols)
health
\end{lstlisting}

\begin{center}
\renewcommand{\arraystretch}{1.2}
\begin{tabular}{l l r r r r}
\toprule
 & \textbf{subject} & \multicolumn{2}{c}{\textbf{Ada}} & \multicolumn{2}{c}{\textbf{Bo}} \\
 & \textbf{type} & \textbf{hr} & \textbf{temp} & \textbf{hr} & \textbf{temp} \\
\textbf{year} & \textbf{visit} & & & & \\
\midrule
2024 & 1 & 74.0 & 37.4 & 78.0 & 36.6 \\
     & 2 & 76.0 & 38.1 & 72.0 & 37.0 \\
2025 & 1 & 80.0 & 36.9 & 74.0 & 37.7 \\
     & 2 & 70.0 & 37.2 & 82.0 & 36.8 \\
\bottomrule
\end{tabular}
\end{center}

Because the columns are hierarchical, a top-level key selects a whole sub-table --- everything for one subject:

\begin{lstlisting}
health['Ada']            # a DataFrame with columns ['hr', 'temp']
health['Ada', 'hr']      # a single Series: Ada's heart rate over time
\end{lstlisting}

\begin{keyidea}
The same partial-indexing rule applies on both axes. A single bare key reaches into the \emph{outermost} level and returns the corresponding lower-dimensional slice. \code{health['Ada']} is to columns what \code{pop['California']} is to rows.
\end{keyidea}

\section{Indexing and slicing multiply-indexed objects}

\subsection{On a Series}

Take the named \code{pop} from above. The pattern is: supply labels for levels left-to-right, and Pandas returns whatever is left.

\begin{lstlisting}
pop['California']          # partial index, outer level -> a Series over year
pop['California', 2010]    # both levels -> a single scalar, 37253956
pop['California':'New York']        # outer-level slice (endpoint INCLUSIVE)
pop[:, 2020]               # fix the INNER level, keep all outer labels
pop[pop > 22_000_000]      # boolean mask works as usual
pop[['California', 'Texas']]        # fancy indexing on the outer level
\end{lstlisting}

\begin{pitfall}
\code{pop[:, 2020]} --- using the empty slice \code{:} as a placeholder to ``keep everything at the outer level while pinning the inner level'' --- only works because this is a \code{Series}. On a \code{DataFrame}, a bare bracket lookup \code{df[...]} indexes \emph{columns}, so this trick does not transfer; you need \code{.loc} and (often) \code{pd.IndexSlice}, covered next.
\end{pitfall}

\subsection{On a DataFrame}

\code{DataFrame} selection goes through \code{.loc} and \code{.iloc}. With a row \code{MultiIndex}, the \code{.loc} key is itself a tuple of per-level labels:

\begin{lstlisting}
health.loc[2024]               # partial: all visits in 2024 (a DataFrame)
health.loc[(2024, 1)]          # both row levels -> one row (a Series)
health.loc[(2024, 1), ('Ada', 'hr')]   # one cell, addressing all 4 levels
\end{lstlisting}

The trouble starts when you want to slice the \emph{inner} level. You might reach for \code{health.loc[(:, 1), :]} to mean ``visit 1 in every year,'' but a bare colon is illegal \emph{inside} a tuple --- Python parses \code{:} as slice syntax only in subscripts, not inside parentheses. That syntax error is what \code{pd.IndexSlice} exists to dodge.

\subsection{\code{pd.IndexSlice} for inner-level slicing}

\code{pd.IndexSlice} is a tiny helper object whose only job is to let you write slice syntax (\code{:}, \code{start:stop}) \emph{inside} a \code{.loc} key, per level:

\begin{lstlisting}
idx = pd.IndexSlice

# every year, but only visit 1; and only the 'hr' column of each subject
health.loc[idx[:, 1], idx[:, 'hr']]
\end{lstlisting}

\begin{outblock}
subject         Ada    Bo
                 hr    hr
year visit
2024 1         74.0  78.0
2025 1         80.0  74.0
\end{outblock}

Read \code{idx[:, 1]} as ``all of level 0 (year), level 1 (visit) equal to 1,'' and \code{idx[:, 'hr']} as the analogous slice across the column levels. \code{IndexSlice} is purely syntactic sugar --- it builds the right tuple of slice objects so you can express row \emph{and} column inner-level selection in a single \code{.loc} call.

\section{Rearranging multi-indexed data}

\subsection{Sorting: why slicing demands a sorted index}

A \code{MultiIndex} stores its labels as \emph{codes} pointing into per-level label arrays. Range slicing on those levels (\code{'California':'New York'}, or any partial outer slice) is only well-defined and only \emph{fast} when the index is \textbf{lexically sorted} --- because a slice is internally a search for a contiguous block of codes, and ``contiguous'' only means anything if the labels are in order. If you try to slice an unsorted \code{MultiIndex}, Pandas refuses:

\begin{lstlisting}
bad = pd.Series(range(4),
                index=pd.MultiIndex.from_tuples(
                    [('b', 1), ('a', 1), ('b', 2), ('a', 2)]))
bad['a':'b']
# UnsortedIndexError: 'Key length (1) was greater than MultiIndex
#   lexsort depth (0)'
\end{lstlisting}

The cure is one call:

\begin{lstlisting}
good = bad.sort_index()
good['a':'b']          # now works -- the block ('a',1),('a',2) is contiguous
\end{lstlisting}

\begin{pitfall}
\code{UnsortedIndexError} (or the older \code{KeyError} about lexsort depth) almost always means ``you sliced a \code{MultiIndex} you never sorted.'' The fix is essentially always \code{df = df.sort\_index()} right after you build or load the data. Sorting once up front makes every later partial slice both legal and fast.
\end{pitfall}

\subsection{\code{reset\_index} and \code{set\_index}}

Sometimes a flat table with the levels as ordinary columns is more convenient (for merging, plotting, or writing to CSV). \code{reset\_index} turns index levels into columns; \code{set\_index} does the reverse.

\begin{lstlisting}
flat = pop.reset_index(name='population')
#   state        year   population
# 0 California   2010    37253956
# 1 California   2020    39538223
# ...
flat.set_index(['state', 'year'])   # back to a MultiIndexed object
\end{lstlisting}

\begin{intuition}
\code{reset\_index} $\leftrightarrow$ \code{set\_index} is the ``index $\leftrightarrow$ column'' round trip, just as \code{stack} $\leftrightarrow$ \code{unstack} is the ``rows $\leftrightarrow$ columns'' round trip. Between these two pairs you can move any label axis to wherever it is most convenient for the operation at hand.
\end{intuition}

\subsection{\code{stack}/\code{unstack} with a chosen level}

When more than two levels are in play, tell \code{unstack} \emph{which} level to pivot using \code{level=} (by position or by name):

\begin{lstlisting}
pop.unstack(level=0)    # 'state' becomes the columns
pop.unstack(level=1)    # 'year' becomes the columns (the default: innermost)
pop.unstack(level='year')   # same, by name -- clearer
\end{lstlisting}

\section{Aggregating across a level}

The real payoff of named levels is group-wise reduction. To collapse one axis --- say, average over years to get a per-state mean --- group by the levels you want to \emph{keep} and reduce:

\begin{lstlisting}
pop.groupby(level='state').mean()
\end{lstlisting}

\begin{outblock}
state
California    38396089.5
New York      19789675.5
Texas         27145533.0
dtype: float64
\end{outblock}

For a \code{DataFrame} with hierarchical columns, \code{groupby} also accepts an \code{axis} and a column-level name, so you can aggregate across, say, subjects while keeping measurement type:

\begin{lstlisting}
# mean over both subjects, per (year, visit) row and per measurement type
health.groupby(level='type', axis=1).mean()
\end{lstlisting}

\begin{pitfall}
You may see older code written as \code{pop.mean(level='state')} --- passing \code{level=} straight to a reduction like \code{mean}, \code{sum}, or \code{std}. That form is \textbf{deprecated} (removed in pandas 2.0). Always spell it \code{pop.groupby(level='state').mean()} instead. It is more explicit about what is happening (group, then reduce) and it is the form that will keep working.
\end{pitfall}

\begin{keyidea}
\code{groupby(level=...)} is just \code{groupby} where the grouping keys come from \emph{index levels} rather than columns. Everything you know about \code{groupby} --- multiple aggregations via \code{.agg}, custom functions, transform --- applies unchanged. The \code{MultiIndex} simply hands \code{groupby} a ready-made set of keys.
\end{keyidea}

\section*{Section recap}
\addcontentsline{toc}{section}{Section recap}
\begin{itemize}
  \item A \code{MultiIndex} packs $N$-dimensional labeled data into a 1-D \code{Series} or 2-D \code{DataFrame} by giving the index multiple \emph{levels}. Tuples-as-keys look similar but lack level structure --- promote them with \code{pd.MultiIndex}.
  \item Build one implicitly (list-of-lists index, or dict with tuple keys) or explicitly via \code{from\_tuples}, \code{from\_arrays}, \code{from\_product}; always set \code{index.names}.
  \item A row \code{MultiIndex} and the column axis are two views of one cube: \code{unstack} widens (level $\to$ columns), \code{stack} narrows (columns $\to$ level). Columns can be hierarchical too.
  \item Index left-to-right by level; a bare outer key returns a lower-dimensional slice. Use \code{.loc} with \code{pd.IndexSlice} to slice \emph{inner} levels on a \code{DataFrame}.
  \item \code{sort\_index()} before slicing a \code{MultiIndex}, or you get \code{UnsortedIndexError}. \code{reset\_index}/\code{set\_index} move levels to and from columns.
  \item Aggregate across a level with \code{groupby(level=...)}; the older \code{.mean(level=...)} form is deprecated.
\end{itemize}

\end{document}
